\begin{table}[!htbp]
\centering
\caption[Comparație runner-e compartimentale]{Comparație side-by-side a celor trei runner-e compartimentale, ca
instanțe ale aceluiași contract de simulare. Diferențele țin de ecosistemul
numeric (biblioteca solver dominantă, convenția de control al pasului), nu de
modelul matematic implementat. Numărul de linii include parsarea mesajului,
dispecerizarea pe model, scrierea în baza de date și colectarea metricilor.}
\label{tab:implementari_runners}
\small
\begin{tabularx}{\columnwidth}{l X l r}
\toprule
\textbf{Limbaj} & \textbf{Bibliotecă \& metodă} & \textbf{Pas} & \textbf{Linii} \\
\midrule
\texttt{r}      & \texttt{deSolve::ode}, RK4 clasic                                  & fix     & 274 \\
\texttt{julia}  & \texttt{DiffEq.jl}, \texttt{Tsit5} (Tsitouras 5(4))                & adaptiv & 236 \\
\texttt{python} & \texttt{scipy.solve\_ivp}, \texttt{RK45} (Dormand-Prince)          & adaptiv & 259 \\
\bottomrule
\end{tabularx}
\end{table}
